\documentclass[../readings.tex]{subfiles}

\begin{document}
\subsection{Numerical Quadrature}
\label{subsec:quadrature}

Approximating $I = \int_a^b f(x)\,dx$ via weighted sums: $\sum_{i=1}^n w_i f(x_i)$.

\subsubsection{Newton-Cotes Rules}

Uses equally-spaced nodes. Higher-order rules are prone to Runge's phenomenon.

\addDef{Trapezoidal Rule}{%
On $n$ panels of width $h = (b-a)/n$:
\[
    T_n = h \left[ \frac{f(a) + f(b)}{2} + \sum_{i=1}^{n-1} f(a + ih) \right].
\]
\textbf{Error Bound:}\enspace $E = O(h^2)$. Doubling $n$ reduces error $4\times$.
}

\addDef{Simpson's Rule}{%
Requires $n$ to be even. Uses piecewise quadratic interpolation:
\[
    S_n = \frac{h}{3} \left[ f(a) + f(b) + 4\sum_{\text{odd}\;i} f(x_i) + 2\sum_{\text{even}\;i} f(x_i) \right].
\]
\textbf{Error Bound:}\enspace $E = O(h^4)$. Exact for polynomials of degree $\leq 3$.
}

\subsubsection{Gaussian Quadrature}

Optimizes \textbf{both} nodes and weights to achieve the highest possible degree of exactness.

\addThm{Gauss-Legendre Quadrature}{%
Integrating over $[-1, 1]$ with $p$ points:
\[
    \int_{-1}^1 f(x)\,dx \approx \sum_{i=1}^p w_i f(x_i).
\]
\textbf{Optimality:}\enspace Exact for all polynomials of degree $\leq 2p-1$.
\textbf{Nodes:}\enspace Roots of the degree-$p$ Legendre polynomial $P_p(x)$.
}

\addNote{%
\textbf{Efficiency:}\enspace $p$ Gauss points are as accurate as $2p$ Newton-Cotes points. Always prefer Gaussian quadrature for smooth, non-periodic functions.
}

\subsubsection{Richardson Extrapolation and Romberg Integration}

\addThm{Richardson Extrapolation}{%
Given an $O(h^2)$ estimate $A(h)$, build an $O(h^4)$ estimate by:
\[
    A_{new} = \frac{4A(h/2) - A(h)}{3}.
\]
}

\addDef{Romberg Integration}{%
Iterative application of Richardson extrapolation to the trapezoidal rule.
\textbf{Performance:}\enspace Error decreases as $O(h^{2k+2})$ for smooth $f$, converging spectacularly fast.
}

\subsubsection{Adaptive Quadrature}

\addNote{%
\textbf{Design Goal:}\enspace Automatically focus function evaluations where $f$ is "difficult" (steep gradients, oscillations) and use large steps where $f$ is smooth.
\textbf{Mechanism:}\enspace Compare quadrature on one large panel vs.\ two small panels. If the difference is above tolerance $\varepsilon$, recurse.
}

\subsubsection{Periodic Functions and Spectral Accuracy}

\addNote{%
\textbf{The Periodic Surprise:}\enspace For smooth, periodic functions on $[0, 2\pi]$, the simple **Trapezoidal Rule** is the optimal choice. It achieves **Spectral Accuracy** (exponential convergence) because error terms in the Euler-Maclaurin expansion vanish at endpoints.
}

\subsubsection{Exercises}

\addExcr{%
\begin{enumerate}
    \item Doubling nodes in Trapezoidal rule reduces error by $4\times$. How much does Simpson reduce it?
    \item Derive the $p=2$ Gauss points for $[-1, 1]$. (\textit{Hint: roots of $3x^2 - 1$}).
    \item Use Richardson extrapolation to improve $T(h)$ and $T(h/2)$ to get Simpson's rule.
    \item Integration of $\sin(x)$ on $[0, 2\pi]$ with $n=4$ nodes. Is the error $0$? Why?
\end{enumerate}
}

\end{document}
